\input{../tex-inputs/latex-korrekturansicht-vorspann}

\section[Gustav Schwarzkopf an Arthur Schnitzler, {[}14. 12. 1905–18. 12. 1905?{]}]{L04239 Gustav Schwarzkopf an Arthur Schnitzler, {[}14. 12. 1905–18. 12. 1905?{]}}
\nopagebreak\mylabel{L04239v}
\rehead{ }\normalsize\beginnumbering\briefempfaengerindex{Schnitzler, Arthur@\textsc{Schnitzler, Arthur}!zzzSchwarzkopf, Gustav@\emph{von Gustav Schwarzkopf}!1905-12-181@{{[}14. 12. 1905–18. 12. 1905?{]}}|(be}
\toendnotes[C]{\smallbreak\pagebreak[2]}
\correspDesc{Versand  durch Gustav Schwarzkopf im Zeitraum [14. 12. 1905 –
                  18. 12. 1905?] in Wien
\newline{}Erhalt  durch Arthur Schnitzler am [18. 12. 1905?] in Wien}\toendnotes[C]{\smallbreak}
\Standort{CUL, Schnitzler, B 96a.}
\physDesc{Brief, 1 Blatt, 4 Seiten, 1246 Zeichen
\newline{}Handschrift: schwarze Tinte, deutsche Kurrent
\newline{}Schnitzler: mit Bleistift datiert: »18/12 05(?)« }\toendnotes[C]{\smallbreak}
\pstart
           \noindent{}{\pb}\label{K_L04239-1v}\edtext{S. 115}{\lemma{\textnormal{\emph{S. 115}}}\Cendnote{\textnormal{Zieht man
            die unsichere Datierung \textcolor{blue}{Schnitzlers} heran, dürfte er bei seinem Besuch bei \textcolor{blue}{Schwarzkopf}\pwindex{Schwarzkopf, Gustav 7.\,11.\,1853 Wien – 13.\,11.\,1939 ebd.@\textsc{Schwarzkopf, Gustav} (7.\,11.\,1853 Wien – 13.\,11.\,1939 ebd.), \emph{Schriftsteller}|pwk}
                  am 14. 12. 1905 ein getipptes Werkmanuskript von \emph{\textcolor{green}{Der Ruf
                     des Lebens}\pwindex{Schnitzler, Arthur 15. 5. 1862 Wien – 21. 10. 1931 ebd.@\textsc{Schnitzler, Arthur} (15. 5. 1862 Wien – 21. 10. 1931 ebd.), \emph{Schriftsteller, Mediziner}!Ruf des Lebens. Schauspiel in drei Akten@\strich\emph{Der Ruf des Lebens. Schauspiel in drei Akten}|pwk}} zurückgelassen haben. Die Seitenzahlen beziehen sich auf keine überlieferte Fassung.
            }}}\label{K_L04239-1}. Der \textcolor{green}{Arzt}\pwindex{Schnitzler, Arthur 15. 5. 1862 Wien – 21. 10. 1931 ebd.@\textsc{Schnitzler, Arthur} (15. 5. 1862 Wien – 21. 10. 1931 ebd.), \emph{Schriftsteller, Mediziner}!Ruf des Lebens. Schauspiel in drei Akten@\strich\emph{Der Ruf des Lebens. Schauspiel in drei Akten}|pwv}{}\ledrightnote{{$\rightarrow$}\emph{\textcolor{green}{Der Ruf des Lebens. Schauspiel in drei Akten}}}. Zu \textcolor{gray}{verſchwo{\geminationm}en
               fü}r die Situation. We{\geminationn} nichts Beſſeres zu
               finden, lieber die ganze Rede weglaſſen und gleich den \textcolor{green}{Adjunkten}\pwindex{Schnitzler, Arthur 15. 5. 1862 Wien – 21. 10. 1931 ebd.@\textsc{Schnitzler, Arthur} (15. 5. 1862 Wien – 21. 10. 1931 ebd.), \emph{Schriftsteller, Mediziner}!Ruf des Lebens. Schauspiel in drei Akten@\strich\emph{Der Ruf des Lebens. Schauspiel in drei Akten}|pwv}{}\ledrightnote{{$\rightarrow$}\emph{\textcolor{green}{Der Ruf des Lebens. Schauspiel in drei Akten}}} auftreten.\pend
           
\pstart
           S. 129. Letze Rede d\textcolor{gray}{es}\textcolor{green}{Adjunkten}\pwindex{Schnitzler, Arthur 15. 5. 1862 Wien – 21. 10. 1931 ebd.@\textsc{Schnitzler, Arthur} (15. 5. 1862 Wien – 21. 10. 1931 ebd.), \emph{Schriftsteller, Mediziner}!Ruf des Lebens. Schauspiel in drei Akten@\strich\emph{Der Ruf des Lebens. Schauspiel in drei Akten}|pwv}{}\ledrightnote{{$\rightarrow$}\emph{\textcolor{green}{Der Ruf des Lebens. Schauspiel in drei Akten}}} die »reine Stirne«
               ſcheint mir gefährlich. Er ſagt auch zu oft »Leben Sie wol“ und »Auf Wiederſehen«
               ſollte er gar nicht ſagen. Vielleicht ungefähr ſo:\pend
           
\pstart
           Der \textcolor{green}{Adjunkt}\pwindex{Schnitzler, Arthur 15. 5. 1862 Wien – 21. 10. 1931 ebd.@\textsc{Schnitzler, Arthur} (15. 5. 1862 Wien – 21. 10. 1931 ebd.), \emph{Schriftsteller, Mediziner}!Ruf des Lebens. Schauspiel in drei Akten@\strich\emph{Der Ruf des Lebens. Schauspiel in drei Akten}|pwv}{}\ledrightnote{{$\rightarrow$}\emph{\textcolor{green}{Der Ruf des Lebens. Schauspiel in drei Akten}}} reicht dem \textcolor{green}{Arzt}\pwindex{Schnitzler, Arthur 15. 5. 1862 Wien – 21. 10. 1931 ebd.@\textsc{Schnitzler, Arthur} (15. 5. 1862 Wien – 21. 10. 1931 ebd.), \emph{Schriftsteller, Mediziner}!Ruf des Lebens. Schauspiel in drei Akten@\strich\emph{Der Ruf des Lebens. Schauspiel in drei Akten}|pwv}{}\ledrightnote{{$\rightarrow$}\emph{\textcolor{green}{Der Ruf des Lebens. Schauspiel in drei Akten}}} die Hand, »Leben Sie wol.
                  (da{\geminationn} wendet er {\pb}ſich zu \textcolor{green}{Marie}\pwindex{Schnitzler, Arthur 15. 5. 1862 Wien – 21. 10. 1931 ebd.@\textsc{Schnitzler, Arthur} (15. 5. 1862 Wien – 21. 10. 1931 ebd.), \emph{Schriftsteller, Mediziner}!Ruf des Lebens. Schauspiel in drei Akten@\strich\emph{Der Ruf des Lebens. Schauspiel in drei Akten}|pwv}{}\ledrightnote{{$\rightarrow$}\emph{\textcolor{green}{Der Ruf des Lebens. Schauspiel in drei Akten}}}) \textcolor{green}{Marie}\pwindex{Schnitzler, Arthur 15. 5. 1862 Wien – 21. 10. 1931 ebd.@\textsc{Schnitzler, Arthur} (15. 5. 1862 Wien – 21. 10. 1931 ebd.), \emph{Schriftsteller, Mediziner}!Ruf des Lebens. Schauspiel in drei Akten@\strich\emph{Der Ruf des Lebens. Schauspiel in drei Akten}|pwv}{}\ledrightnote{{$\rightarrow$}\emph{\textcolor{green}{Der Ruf des Lebens. Schauspiel in drei Akten}}}, – ſie lange betrachtend) Wie Eine, die aus tiefſter
               Ferne zurückgekehrt iſt, erſcheinen. Sie mir – (wie ſtaunend) Ihre Sti{\geminationm}e klingt mir wieder wir ſie mir früher einmal geklungen
               – ich empfinde kein Grauen mehr – nur tiefes, i{\geminationn}igs
               Mitleid — (reicht ihr die Hand, die ſie zögernd ni{\geminationm}t und
               hält ſie einige Sekunden, da{\geminationn} mit Blick ab. \pend
           
\pstart
           {\pb}S. 129\pend
           
\pstart
           \uline{der \textcolor{green}{Arzt}\pwindex{Schnitzler, Arthur 15. 5. 1862 Wien – 21. 10. 1931 ebd.@\textsc{Schnitzler, Arthur} (15. 5. 1862 Wien – 21. 10. 1931 ebd.), \emph{Schriftsteller, Mediziner}!Ruf des Lebens. Schauspiel in drei Akten@\strich\emph{Der Ruf des Lebens. Schauspiel in drei Akten}|pwv}{}\ledrightnote{{$\rightarrow$}\emph{\textcolor{green}{Der Ruf des Lebens. Schauspiel in drei Akten}}}} Sie werden ihn wiederſehn \textcolor{green}{Marie}\pwindex{Schnitzler, Arthur 15. 5. 1862 Wien – 21. 10. 1931 ebd.@\textsc{Schnitzler, Arthur} (15. 5. 1862 Wien – 21. 10. 1931 ebd.), \emph{Schriftsteller, Mediziner}!Ruf des Lebens. Schauspiel in drei Akten@\strich\emph{Der Ruf des Lebens. Schauspiel in drei Akten}|pwv}{}\ledrightnote{{$\rightarrow$}\emph{\textcolor{green}{Der Ruf des Lebens. Schauspiel in drei Akten}}}. Was er jetzt weiß, es wird verblaſſen, es wird ſeine letzten Schwur
               für ihn verlieren.\pend
           
\pstart
           Seite 130\pend
           
\pstart
           Ein ſtärkers Betonen ihrer eigentlichen Schuld \textcolor{gray}{empſendung}. Nach,
               „wie darf ich?«– Begreifen Sie es de{\geminationn}? ich habe das
               Abſcheulichſte getan, ich habe das Gräßlichſte erlebt und ich will noch
                  le\textcolor{gray}{b}en, ich weiß, daß ich {\pb}nicht mehr leben darf, und es lo\uline{ck}t mich mich zu leben. Daß es ſolche Minuten gibt
               u. ſ. w.\pend
           
\pstart
           Schlußwort des \textcolor{green}{Arztes}\pwindex{Schnitzler, Arthur 15. 5. 1862 Wien – 21. 10. 1931 ebd.@\textsc{Schnitzler, Arthur} (15. 5. 1862 Wien – 21. 10. 1931 ebd.), \emph{Schriftsteller, Mediziner}!Ruf des Lebens. Schauspiel in drei Akten@\strich\emph{Der Ruf des Lebens. Schauspiel in drei Akten}|pwv}{}\ledrightnote{{$\rightarrow$}\emph{\textcolor{green}{Der Ruf des Lebens. Schauspiel in drei Akten}}}
               klingt ſehr gut.\pend
           
\pstart
           Herzliche Grüße{\\[\baselineskip]}Ihr{\\[\baselineskip]}\spacefill\mbox{GustavSch}\pend
           \leftskip=0em{}\selectlanguage{ngerman}\endnumbering\briefempfaengerindex{Schnitzler, Arthur@\textsc{Schnitzler, Arthur}!zzzSchwarzkopf, Gustav@\emph{von Gustav Schwarzkopf}!1905-12-141@{{[}14. 12. 1905–18. 12. 1905?{]}}|)be}\mylabel{L04239h}
\begin{anhang}
\end{anhang}\newcommand{\dateiname}{L04239}\newcommand{\titel}{Gustav Schwarzkopf an Arthur Schnitzler, [14. 12. 1905 – 18. 12. 1905?]}\newcommand{\editorInnen}{Herausgegeben von Jahnke, SelmaMüller, Martin Anton}\input{../tex-inputs/latex-korrekturansicht-abspann}
